\documentclass[11pt]{article}

\usepackage[utf8]{inputenc}
\usepackage[T1]{fontenc}
\usepackage{amsmath}
\usepackage{amssymb}
\usepackage{geometry}
\usepackage{hyperref}

\geometry{margin=1in}

\title{Monte Carlo Maze Energy Functions}
\author{}
\date{}

\begin{document}

\maketitle

\section{Genel Kural}

MCMC sisteminde her state $s$ icin bir enerji degeri hesaplanir.
Enerji ne kadar dusukse state o kadar iyi kabul edilir.

\[
E(s) \downarrow \quad \Longrightarrow \quad s \text{ daha iyi}
\]

Metropolis-Hastings kabul kuralinda bu enerji farki kullanilir.

\section{Temel Notasyon}

\begin{tabular}{|l|p{0.68\linewidth}|}
\hline
Sembol & Anlam \\
\hline
$s$ & MCMC state: grid, start, door, itemlar, initial stamina ve locked door bilgisi. \\
\hline
$I$ & Initial stamina. Kodda \texttt{state.initial\_stamina}. \\
\hline
$D^\star$ & Target agent difficulty. Kodda \texttt{TARGET\_AGENT\_DIFFICULTY}. Aralik olarak $[0,1]$ varsayilir. \\
\hline
$W_D$ & Difficulty term agirligi. Kodda \texttt{AGENT\_DIFFICULTY\_WEIGHT}. \\
\hline
$W_T$ & Segment target term agirligi. Kodda \texttt{SEGMENT\_TARGET\_WEIGHT}. \\
\hline
$W_B$ & Segment balance term agirligi. Kodda \texttt{SEGMENT\_BALANCE\_WEIGHT}. \\
\hline
$W_H$ & Dead segment term agirligi. Kodda \texttt{DEAD\_SEGMENT\_WEIGHT}. \\
\hline
$W_U$ & Stamina usage term agirligi. Kodda \texttt{STAMINA\_USAGE\_WEIGHT}. \\
\hline
$W_F$ & Final stamina term agirligi. Kodda \texttt{FINAL\_STAMINA\_WEIGHT}. \\
\hline
$W_X$ & Spacing term agirligi. Kodda \texttt{SPACING\_WEIGHT}. \\
\hline
$F$ & Final stamina target factor. Kodda \texttt{FINAL\_STAMINA\_TARGET\_FACTOR}. \\
\hline
$K_X$ & Spacing target scale. Kodda \texttt{SPACING\_TARGET\_SCALE}. \\
\hline
\end{tabular}

\section{\texttt{door\_only} Energy}

Bu en basit baseline enerji fonksiyonudur.
Sadece start ile door arasindaki en kisa yol uzunluguna bakar.

\[
E_{\text{door}}(s)
= \left| L(s) - L^\star \right|
\]

Burada:

\begin{itemize}
  \item $L(s)$: grid uzerinde \texttt{start -> door} en kisa yol uzunlugu.
  \item $L^\star$: hedef yol uzunlugu. Kodda \texttt{TARGET\_PATH\_LENGTH}.
\end{itemize}

Eger door'a ulasilamiyorsa:

\[
E_{\text{door}}(s) = \infty
\]

Bu modda stamina, key ve item semantigi yoktur.

\section{\texttt{stamina\_only} Baseline Energy}

Bu eski stamina-aware baseline enerji fonksiyonudur.
HC3 solver exact olarak haritanin cozulebilir olup olmadigini hesaplar.

\[
E_{\text{base}}(s)
=
w_P \left| P(s) - P^\star \right|
+
w_R \left| R_{\text{best}}(s) - R^\star \right|
\]

Burada:

\begin{itemize}
  \item $P(s)$: exact basarili cozumler icindeki en kisa basarili cozum maliyeti.
  \item $P^\star$: hedef cozum uzunlugu. Kodda \texttt{TARGET\_PATH\_LENGTH}.
  \item $R_{\text{best}}(s)$: door'a varinca kalan en iyi stamina.
  \item $R^\star$: hedef final stamina. Kodda \texttt{TARGET\_FINAL\_STAMINA}.
  \item $w_P$: path term agirligi.
  \item $w_R$: remaining stamina term agirligi.
\end{itemize}

Exact cozum path maliyeti su sekilde yorumlanir:

\[
P(s)
=
I
+
C_{\text{stamina}}
-
R_{\text{final}}
\]

Burada:

\begin{itemize}
  \item $I$: initial stamina.
  \item $C_{\text{stamina}}$: cozum boyunca toplanan stamina item degerlerinin toplami.
  \item $R_{\text{final}}$: door'a ulasildiginda kalan stamina.
\end{itemize}

Harita exact olarak cozulemezse:

\[
E_{\text{base}}(s) = \infty
\]

\section{Agent Difficulty Energy}

Guncel difficulty-aware stamina energy fonksiyonu birden fazla terimden olusur:

\[
E_{\text{agent}}(s)
=
E_D(s)
+
E_T(s)
+
E_B(s)
+
E_H(s)
+
E_U(s)
+
E_F(s)
+
E_X(s)
\]

Terimler:

\begin{itemize}
  \item $E_D$: main route difficulty hedefinden sapma.
  \item $E_T$: segmentlerin hedef lokal success rate'ten sapmasi.
  \item $E_B$: segmentler arasi zorluk dengesizligi.
  \item $E_H$: exact dead segment orani.
  \item $E_U$: basarili planlarin stamina item toplama orani hedefinden sapma.
  \item $E_F$: final stamina hedefinden sapma.
  \item $E_X$: start-key-door spacing hedefinden sapma.
\end{itemize}

\subsection{Semantic Plan ve Agent Success Rate}

HC3 semantic graph uzerinde door'a giden simple semantic planlar uretilir.
Exact olarak cozulebilen planlar kumesi:

\[
\mathcal{P}_{\text{solv}}
=
\{p_i \mid p_i \text{ exact-solvable}\}
\]

Her exact-solvable plan $p_i$ segmentlerden olusur:

\[
p_i = (e_{i,1}, e_{i,2}, \ldots, e_{i,m_i})
\]

Segment $e_{i,j}$ icin agent success rate:

\[
q_{i,j}
=
\frac{
  \#\text{successful agents on segment } e_{i,j}
}{
  \#\text{agents simulated on segment } e_{i,j}
}
\]

Plan success rate segment success rate carpimidir:

\[
R_i
=
\prod_{j=1}^{m_i} q_{i,j}
\]

Exact-dead planlar icin agent calistirilmaz ve plan success rate $0$ kabul edilir.
Ancak bu planlar main route average'a katilmaz; dead segment term icinde temsil edilir.

\subsection{Difficulty Term}

Exact-solvable planlarin ortalama success rate'i:

\[
\bar{R}
=
\frac{1}{|\mathcal{P}_{\text{solv}}|}
\sum_{p_i \in \mathcal{P}_{\text{solv}}} R_i
\]

Main route difficulty:

\[
D_{\text{main}}
=
1 - \bar{R}
\]

Difficulty energy:

\[
E_D(s)
=
W_D
\left|
D_{\text{main}}
-
D^\star
\right|
\]

Yorum:

\begin{itemize}
  \item $D_{\text{main}}$ hedefe yakin ise bu terim dusuk olur.
  \item Agentlar cok kolay geciyorsa $\bar{R}$ yuksek, $D_{\text{main}}$ dusuk olur.
  \item Agentlar zor geciyorsa $\bar{R}$ dusuk, $D_{\text{main}}$ yuksek olur.
\end{itemize}

\subsection{Segment Target Term}

Main route hedefi once route success hedefi olarak yazilir:

\[
R^\star
=
1 - D^\star
\]

Plan $p_i$ toplam $m_i$ segmentten olusuyorsa, o plan icin lokal segment success hedefi:

\[
Q_i^\star
=
(R^\star)^{1/m_i}
\]

Segment target score:

\[
T_{\text{score}}
=
\frac{1}{M}
\sum_{p_i \in \mathcal{P}_{\text{solv}}}
\sum_{j=1}^{m_i}
\left|
q_{i,j} - Q_i^\star
\right|
\]

Burada $M$ exact-solvable planlardaki toplam simulated segment occurrence sayisidir.
Kodda segment yoksa skor $0$ kabul edilir ve skor $1$ degerinde saturate edilir.

Segment target energy:

\[
E_T(s)
=
W_T T_{\text{score}}
\]

Yorum:

\begin{itemize}
  \item Bu terim segmentleri route hedefinden turetilen lokal basari oranina yaklastirir.
  \item Eski std tabanli balance sadece segmentlerin birbirine benzerligini olcerken, bu terim segmentlerin hedef difficulty'ye uygun olup olmadigini da olcer.
  \item Tum segmentler cok kolaysa veya cok zorsa bu terim ceza uretir.
\end{itemize}

\subsection{Segment Balance Term}

Unique simulated segmentler kumesi:

\[
\mathcal{U}
=
\{u_1, u_2, \ldots, u_n\}
\]

Ayni semantic segment birden fazla plan context'inde gorunurse, o segmentin success rate'i ortalama ile tekillestirilir:

\[
a_k
=
\text{average success rate for unique segment } u_k
\]

Bu segment success rate'lerinin ortalamasi:

\[
\bar{a}
=
\frac{1}{n}
\sum_{k=1}^{n} a_k
\]

Standart sapma:

\[
\sigma_{\text{seg}}
=
\sqrt{
  \frac{1}{n}
  \sum_{k=1}^{n}
  (a_k - \bar{a})^2
}
\]

Normalize balance score:

\[
B_{\text{score}}
=
\min(1, 2\sigma_{\text{seg}})
\]

Segment balance energy:

\[
E_B(s)
=
W_B B_{\text{score}}
\]

Yorum:

\begin{itemize}
  \item Segment success rate'leri birbirine yakin ise $\sigma_{\text{seg}}$ dusuk olur.
  \item Bir segment cok kolay, digeri cok zor ise $\sigma_{\text{seg}}$ yuksek olur.
  \item Bu terim zorlugun tek bir segmente yigilmamasini hedefler.
\end{itemize}

\subsection{Dead Segment Term}

Exact-dead planlarda ilk exact basarisiz semantic segment dead segment olarak sayilir.
Unique dead segment kumesi:

\[
\mathcal{H}
\]

Tum ilgili unique segmentler:

\[
\mathcal{V}
=
\mathcal{U}
\cup
\mathcal{H}
\]

Dead segment ratio:

\[
H_{\text{ratio}}
=
\frac{|\mathcal{H}|}{|\mathcal{V}|}
\]

Dead segment energy:

\[
E_H(s)
=
W_H H_{\text{ratio}}
\]

Yorum:

\begin{itemize}
  \item Bu terim exact olarak hic calismayan semantic baglantilari cezalandirir.
  \item Dead route suffix'leri tekrar tekrar sayilmaz; unique dead segment sayilir.
  \item Eger \texttt{DEAD\_SEGMENT\_WEIGHT = 0} ise bu terim enerjiye etki etmez ama summary'de gorunmeye devam eder.
\end{itemize}

\subsection{Stamina Usage Term}

Haritadaki toplam stamina item sayisi:

\[
N_{\text{stam}}
=
\#\text{stamina items on map}
\]

Plan $p_i$ icin toplanan stamina item sayisi:

\[
C_i
=
\#\text{stamina items collected by } p_i
\]

Plan seviyesinde stamina collection ratio:

\[
U_i
=
\begin{cases}
0, & N_{\text{stam}} = 0 \\
\frac{C_i}{N_{\text{stam}}}, & N_{\text{stam}} > 0
\end{cases}
\]

Success-rate agirlikli stamina usage:

\[
\bar{U}
=
\frac{
  \sum_{p_i \in \mathcal{P}_{\text{solv}}}
  R_i U_i
}{
  \sum_{p_i \in \mathcal{P}_{\text{solv}}}
  R_i
}
\]

Kodda payda $0$ ise $\bar{U}=0$ kabul edilir.

Target stamina usage:

\[
U^\star
=
\begin{cases}
0, & \text{haritada stamina item yoksa} \\
\texttt{TARGET\_STAMINA\_USAGE\_RATE}, & \text{explicit target verildiyse} \\
D^\star, & \text{aksi halde}
\end{cases}
\]

Stamina usage score:

\[
U_{\text{score}}
=
\left|
U^\star - \bar{U}
\right|
\]

Stamina usage energy:

\[
E_U(s)
=
W_U U_{\text{score}}
\]

Yorum:

\begin{itemize}
  \item Bu terim stamina item toplamadan biten direct planlar baskin oldugunda ceza uretir.
  \item Haritadaki toplam stamina sayisini hesaba katar: 4 stamina item olan haritada 1 item toplamak $0.25$ usage demektir.
  \item Success-rate agirligi kullanildigi icin ajanlarin pratikte daha cok basardigi planlar daha fazla etki eder.
  \item Daha sert stamina zorunlulugu icin \texttt{TARGET\_STAMINA\_USAGE\_RATE = 0.8} gibi sabit hedef verilebilir.
\end{itemize}

\subsection{Final Stamina Term}

Hedef final stamina:

\[
T_F
=
F \cdot I \cdot (1 - D^\star)
\]

Burada difficulty ile ters oranti vardir.
Zorluk hedefi yuksekse, door'a ulasildiginda daha az stamina kalmasi beklenir.

Plan $p_i$ icin basarili agent final stamina ortalamasi:

\[
S_i
=
\text{average final stamina samples for plan } p_i
\]

Success-rate agirlikli final stamina:

\[
\bar{S}
=
\frac{
  \sum_{p_i \in \mathcal{P}_{\text{solv}}}
  R_i S_i
}{
  \sum_{p_i \in \mathcal{P}_{\text{solv}}}
  R_i
}
\]

Kodda payda $0$ ise $\bar{S}=0$ kabul edilir.

Normalize final stamina score:

\[
F_{\text{score}}
=
\min
\left(
1,
\frac{
  |T_F - \bar{S}|
}{
  \max(1, I)
}
\right)
\]

Final stamina energy:

\[
E_F(s)
=
W_F F_{\text{score}}
\]

Yorum:

\begin{itemize}
  \item Bu terim cozum sonunda kalan stamina'yi hedefler.
  \item Yuksek difficulty icin hedef final stamina dusuktur.
  \item Dusuk difficulty icin hedef final stamina daha yuksektir.
  \item Normalizasyon initial stamina ile yapilir.
\end{itemize}

\subsection{Spacing Term}

Grid scale:

\[
G
=
\sqrt{\texttt{GRID\_WIDTH} \cdot \texttt{GRID\_HEIGHT}}
\]

Key varsa actual spacing:

\[
X_{\text{actual}}
=
\frac{
  d(\text{start}, \text{key})
  +
  d(\text{key}, \text{door})
}{2}
\]

Key yoksa:

\[
X_{\text{actual}}
=
d(\text{start}, \text{door})
\]

Burada $d(a,b)$ grid uzerinde shortest-path mesafesidir.

Spacing target:

\[
X^\star
=
K_X
\cdot
G
\cdot
(0.5 + 0.5D^\star)
\]

Normalize spacing score:

\[
X_{\text{score}}
=
\min
\left(
1,
\left(
  \frac{
    |X^\star - X_{\text{actual}}|
  }{
    \max(1, X^\star)
  }
\right)^2
\right)
\]

Spacing energy:

\[
E_X(s)
=
W_X X_{\text{score}}
\]

Yorum:

\begin{itemize}
  \item Key varsa sadece \texttt{start -> key} ve \texttt{key -> door} bacaklari ortalamaya girer.
  \item Key yoksa \texttt{start -> door} mesafesi kullanilir.
  \item Bu terim ozellikle \texttt{start -> key} ve \texttt{key -> door} bacaklarinin birbirine yapismasini cezalandirir.
  \item Karesel skor hedef spacing'den buyuk sapmalari lineer cezaya gore daha sert buyutur ve $1$ degerinde saturate olur.
  \item Difficulty arttikca target spacing artar.
  \item Mesafe Manhattan degil, shortest-path mesafesidir.
\end{itemize}

\section{Auto Initial Stamina Formula}

\texttt{stamina\_only} config'inde \texttt{initial\_stamina = None} ise initial stamina otomatik uretilir.
Bu deger MCMC proposal'lari sirasinda degismez.

Grid scale:

\[
G
=
\sqrt{\texttt{GRID\_WIDTH} \cdot \texttt{GRID\_HEIGHT}}
\]

Ortalama initial stamina:

\[
\mu_I
=
G
\left(
B_I
+
A_I D^\star
\right)
\]

Burada:

\begin{itemize}
  \item $B_I$: \texttt{INITIAL\_STAMINA\_BASE\_SCALE}
  \item $A_I$: \texttt{INITIAL\_STAMINA\_DIFFICULTY\_SCALE}
\end{itemize}

Normal sapma:

\[
\sigma_I
=
G
\cdot
\texttt{INITIAL\_STAMINA\_NOISE\_STD\_SCALE}
\]

Sample:

\[
I_{\text{sample}}
\sim
\mathcal{N}(\mu_I, \sigma_I)
\]

Clamp limitleri:

\[
I_{\min}
=
G
\cdot
\texttt{MIN\_INITIAL\_STAMINA\_SCALE}
\]

\[
I_{\max}
=
G
\cdot
\texttt{MAX\_INITIAL\_STAMINA\_SCALE}
\]

Son initial stamina:

\[
I
=
\operatorname{round}
\left(
\operatorname{clamp}
\left(
I_{\text{sample}},
I_{\min},
I_{\max}
\right)
\right)
\]

Yorum:

\begin{itemize}
  \item Difficulty arttikca ortalama initial stamina artar.
  \item Random sapma seed'e baglidir.
  \item Ayni seed ayni initial stamina degerini uretir.
  \item Bu deger state proposal degiskeni degildir; sadece initial state generation sirasinda secilir.
\end{itemize}

\section{Terminal Breakdown Alanlari}

\begin{tabular}{|l|p{0.66\linewidth}|}
\hline
Terminal alani & Formuldeki karsiligi \\
\hline
\texttt{main\_route\_difficulty} & $D_{\text{main}}$ \\
\hline
\texttt{main\_route\_success\_rate} & $\bar{R}$ \\
\hline
\texttt{main\_route\_term} & $E_D$ \\
\hline
\texttt{segment\_success\_target} & average $Q_i^\star$ \\
\hline
\texttt{segment\_target\_score} & $T_{\text{score}}$ \\
\hline
\texttt{segment\_target\_term} & $E_T$ \\
\hline
\texttt{segment\_success\_std} & $\sigma_{\text{seg}}$ \\
\hline
\texttt{segment\_balance\_score} & $B_{\text{score}}$ \\
\hline
\texttt{segment\_balance\_term} & $E_B$ \\
\hline
\texttt{dead\_segment\_ratio} & $H_{\text{ratio}}$ \\
\hline
\texttt{dead\_segment\_term} & $E_H$ \\
\hline
\texttt{stamina\_usage\_target} & $U^\star$ \\
\hline
\texttt{stamina\_usage\_actual} & $\bar{U}$ \\
\hline
\texttt{stamina\_usage\_score} & $U_{\text{score}}$ \\
\hline
\texttt{stamina\_usage\_term} & $E_U$ \\
\hline
\texttt{target\_weighted\_final\_stamina} & $T_F$ \\
\hline
\texttt{weighted\_final\_stamina} & $\bar{S}$ \\
\hline
\texttt{final\_stamina\_score} & $F_{\text{score}}$ \\
\hline
\texttt{final\_stamina\_term} & $E_F$ \\
\hline
\texttt{spacing\_target} & $X^\star$ \\
\hline
\texttt{spacing\_actual} & $X_{\text{actual}}$ \\
\hline
\texttt{spacing\_score} & $X_{\text{score}}$ \\
\hline
\texttt{spacing\_term} & $E_X$ \\
\hline
\texttt{total} & $E_{\text{agent}}$ \\
\hline
\end{tabular}

\end{document}
